% theme: biology_and_cells
\MajorQuestion{生物と細胞}
\SetRowSpaceQanda

\Instruction{次の問いに答えなさい。}
\QQ{観察するものをのせるガラス板を何というか。}{スライドガラス}
\QQ{生物の体をつくる、小さな部屋のようなつくりを\NamedBlank{cell}という。}{細胞}
\QQ{観察するものに近い側のレンズを\NamedBlank{objective}という。}{対物レンズ}
\QQ{植物細胞にあり、光合成を行う部分を何というか。}{葉緑体}
\QQ{観察するものの上にかぶせる薄いガラスを\NamedBlank{coverglass}という。}{カバーガラス}
\QQ{目を近づけてのぞくレンズを\NamedBlank{eyepiece}という。}{接眼レンズ}
\QQ{観察するものに水を落とすときに使う器具を何というか。}{スポイト}
\QQ{\RefBlank{cell}の中にある、丸い形をした部分を何というか。}{核}
\QQ{プレパラートをのせる台を何というか。}{ステージ}
\QQ{\RefBlank{objective}を切り替えるために回す部分を何というか。}{レボルバー}
\QQ{体が1つの\RefBlank{cell}だけでできている生物を何というか。}{単細胞生物}
\QQ{\RefBlank{eyepiece}が10倍、\RefBlank{objective}が40倍のとき、顕微鏡の倍率は何倍か。}{400倍}
\QQ{植物細胞にあり、内部に液体を含む部分を何というか。}{液胞}
\QQ{\RefBlank{coverglass}をかぶせるときに使う器具を何というか。}{柄つき針}
\QQ{\RefBlank{cell}が酸素を取り入れ、二酸化炭素を出すはたらきを何というか。}{細胞の呼吸}
\QQ{植物細胞で、細胞膜の外側にある丈夫なつくりを何というか。}{細胞壁}
\QQ{\RefBlank{cell}の中で、核のまわりにある部分を何というか。}{細胞質}
\QQ{\RefBlank{coverglass}をかぶせるとき、何が入らないように注意するか。}{気泡（空気の泡）}
\QQ{\RefBlank{cell}のうち、内側の物質の外側を包む薄い膜を何というか。}{細胞膜}
\QQ{体が多くの\RefBlank{cell}でできている生物を何というか。}{多細胞生物}
